\subsection{Shao--Wu-Inspired Nonlinear Generation}
\label{sec:implemented-shao-wu}

The generator offers two standard positions for the hyperbolic tangent:
\texttt{product\_lags} applies it to the lag factors before multiplication,
whereas \texttt{product\_output} applies it once after the complete product.
For two historical inputs $z_1,z_2$, these choices give
$\tanh(z_1)\tanh(z_2)$ and $\tanh(z_1z_2)$, respectively. A term containing
only a coefficient and one lag remains linear in both modes. The selected
mode is retained with the mechanism. Each mode carries a verification record
stating the particular property checked below.

\subsubsection{Desired Property and Shao--Wu's Theorem}
\label{sec:implemented-shao-wu-theory}
A generated series should retain useful stochastic variation while its
magnitudes remain controlled. We also seek a stronger property: the effect
of an arbitrarily supplied initial history should disappear. Imagine running
the same equation twice from different histories, then supplying exactly the
same sampled noise value to both runs at every future step. A decreasing
distance then measures the loss of dependence on their initial histories.

\citet[Section~3, Eq.~(3.1), electronic reprint p.~5]{shao2007spectral}
formalize this comparison as \emph{geometric moment contraction} (GMC).
For the second-moment version used here,
\begin{equation}
 \mathbb E\|\mathbf x_{\tau+v}-\mathbf x'_{\tau+v}\|_2^2
 \leq C_{\mathrm{GMC}}\rho_{\mathrm{GMC}}^{v},
 \qquad C_{\mathrm{GMC}}>0,\quad 0<\rho_{\mathrm{GMC}}<1.
 \label{eq:implemented-gmc-forgetting}
\end{equation}
The expectation averages squared distances across repeated noise
realizations. The upper bound decreases geometrically; individual
distances can fluctuate from step to step. In the paper's stationary
comparison the histories arise from independent past innovation sequences,
and the future innovation values are shared. Theorem~5.1 also supplies
forgetting bounds for fixed initial histories in its proof.

\paragraph{Theorem used for the factor mode.}
Let $\mathbf{s}_t=(\mathbf x_t,\ldots,\mathbf x_{t-L+1})$ store the $L$
historical vectors required for the next update, and write
\begin{equation}
 \mathbf x_{t+1}=\mathcal R_{\mathcal T}(\mathbf{s}_t;
                                      \boldsymbol\varepsilon_{t+1})
 =F_{\mathcal T}(\mathbf{s}_t)+\boldsymbol\varepsilon_{t+1}.
 \label{eq:implemented-gmc-recurrence}
\end{equation}
Here $F_{\mathcal T}$ is the deterministic expression executed by the AST.
We use the following specialization of
\citet[Theorem~5.1 and its proof, electronic reprint pp.~11--12]{shao2007spectral}
to second moments:
\begin{enumerate}
\item The update rule is measurable and fixed over time. Its innovations
are independent and identically distributed over time.
\item The next observation has a finite second moment at some fixed
reference history. We use the zero history:
$\|\mathcal R_{\mathcal T}(\mathbf0;\boldsymbol\varepsilon)\|_{L^2}<\infty$,
where $\|Z\|_{L^2}=(\mathbb E\|Z\|_2^2)^{1/2}$.
\item For every pair of possible histories, the same-innovation comparison
satisfies
\begin{equation}
\begin{split}
 \|\mathcal R_{\mathcal T}(\mathbf{s};\boldsymbol\varepsilon)
      -\mathcal R_{\mathcal T}(\mathbf{s}';\boldsymbol\varepsilon)\|_{L^2}
 &\leq\sum_{d=0}^{L-1}a_d^{\mathrm{GMC}}
                       \|\mathbf{s}_d-\mathbf{s}'_d\|_2,\\
 a_d^{\mathrm{GMC}}&\geq0,\qquad
 c_{\mathrm{GMC}}:=\sum_{d=0}^{L-1}a_d^{\mathrm{GMC}}<1.
\end{split}
\label{eq:implemented-gmc-condition}
\end{equation}
The vector $\mathbf{s}_d\in\mathbb R^N$ contains all channels at offset $d$.
Each $a_d^{\mathrm{GMC}}$ is a sensitivity budget for one historical vector;
their sum limits the combined feedback through all channels and lags.
\end{enumerate}
These conditions yield a causal strictly stationary solution with GMC:
present and past innovations determine the solution, and its joint
distributions are unchanged by time shifts. The finite histories supplied
to the simulator have their own declared distribution, so early generated
observations can show an initialization transient.

\subsubsection{Factor Saturation and Contraction Verification}
The implementation in \path{src/tsfmxai/generators/shao_wu.py}
samples a tree of sums and products with finite coefficients and positive
lags. It then performs the following checks.

\subparagraph{1. Saturate the factors that create nonlinear feedback.}
When a product combines two lag-dependent factors, its lag leaves receive
$\tanh$ before the product is evaluated. For example,
$a x_t x_{t-1}+b x_{t-2}$ becomes
$a\tanh(x_t)\tanh(x_{t-1})+b x_{t-2}$.
This controls the magnitudes of the nonlinear factors and their derivatives.
The last term keeps its linear dependence; its derivative is the finite
coefficient $b$. Every output uses fixed parameters and i.i.d. Gaussian
innovations with covariance $\sigma_\varepsilon^2 I$.

\subparagraph{2. Calculate conservative bounds along the AST.}
For each subtree $e$, let $A(e)$ bound its absolute value and let $D(e)$
contain nonnegative absolute derivative bounds by source and lag.
\begin{table}[H]
\centering\small
\begin{tabular}{@{}lll@{}}
\toprule
Subexpression & Amplitude bound $A$ & Derivative bounds $D$\\
\midrule
Constant $c$ & $|c|$ & $0$\\
Lag leaf & $+\infty$ & $1$ at its source/lag, $0$ elsewhere\\
$\tanh(e)$ & $1$ & $D(e)$\\
$e_1+e_2$ & $A(e_1)+A(e_2)$ & $D(e_1)+D(e_2)$\\
$e_1e_2$ & $A(e_1)A(e_2)$ & $A(e_2)D(e_1)+A(e_1)D(e_2)$\\
\bottomrule
\end{tabular}
\caption{Project rules for global size and derivative bounds. A zero
derivative contributes zero even when the other amplitude bound is infinite.
The $\tanh$ row uses $|\tanh'(z)|\leq1$ and can overestimate the true derivative.}
\label{tab:ast-global-bounds}
\end{table}
For a product, a perturbation in either factor is multiplied by the other
factor's magnitude; this is the product rule. Saturating the nonlinear
factors makes these upper bounds finite. A free linear term can have
unbounded amplitude and a finite derivative. The theorem uses the
derivative budgets and the moment at the reference history, so such terms
can remain in the factor-mode construction.

\subparagraph{3. Scale the combined sensitivity below one.}
For offset $d$, stack all target-by-source derivative bounds into an
$N\times N$ nonnegative matrix $\mathbf D_d$. Its largest singular value
bounds the amplification of a change in that whole historical vector.
Set
\begin{equation}
 \widetilde a_d=\|\mathbf D_d\|_2,\qquad
 \widetilde c=\sum_d\widetilde a_d,\qquad
 \eta_{\mathrm{GMC}}=
 \eta_{\mathrm{safe}}\min\{1,g_\star/\widetilde c\}.
 \label{eq:implemented-gmc-scaling}
\end{equation}
When $\widetilde c=0$, the minimum is defined as one. The defaults
$g_\star=0.8$ and $\eta_{\mathrm{safe}}=0.95$ give a final summed budget
$c_{\mathrm{GMC}}\leq0.76$. All output trees receive the same calculated
root multiplier, stored as \texttt{GMC\_SCALE}. It is computed after
sampling the structure and bindings. This changes the executed
coefficients while preserving the lag support.

Integrating the derivative bound along the line between two histories
gives Eq.~\eqref{eq:implemented-gmc-condition}. The additive innovation
cancels in that comparison because its numerical value is shared.
Using a whole matrix for each lag accounts for feedback between channels.

\subparagraph{4. Check the reference moment and retain the verification record.}
With the scaled drift $F_{\mathcal T}$, the Gaussian innovation law gives
\begin{equation}
 \|\mathcal R_{\mathcal T}(\mathbf0;\boldsymbol\varepsilon)\|_{L^2}
 =\left(\|F_{\mathcal T}(\mathbf0)\|_2^2+
                   N\sigma_\varepsilon^2\right)^{1/2}<\infty.
 \label{eq:implemented-gmc-anchor}
\end{equation}
Thus the theorem's reference-moment condition is checked directly even
when the drift is unbounded away from zero.
\texttt{ShaoWuCertificate} stores the lag budgets, their sum, the reference
moment, and the innovation moment. Its guarantee is identified as
\path{geometric_moment_contraction}.

For example, the proposed drift $0.3\tanh(x_t)\tanh(x_{t-1})$ has two
derivative bounds of $0.3$. Applying the default safety factor gives
\begin{equation}
 x_{t+1}=0.285\tanh(x_t)\tanh(x_{t-1})+\varepsilon_{t+1},
 \qquad \varepsilon_{t+1}\sim\mathcal N(0,0.05^2).
 \label{eq:implemented-gmc-example}
\end{equation}
The final total sensitivity bound is $0.57<1$ and the reference moment is
finite. This is a worked project construction using the cited theorem.

\subsubsection{Product-Output Saturation and Moment Bounds}
The second standard option evaluates the product first, then applies
$\tanh$ once. Its construction keeps history-independent multiplicative
coefficients outside the saturation: both $a(z_1z_2)$ and $(az_1)z_2$
therefore produce $a\tanh(z_1z_2)$. For three or more factors the entire
connected product is saturated once. Additive terms keep their original
roles.

This choice directly bounds the product's output. Its sensitivity has a
different behaviour:
\begin{equation}
 \frac{\partial\,\tanh(z_1z_2)}{\partial z_1}
 =z_2\bigl[1-\tanh^2(z_1z_2)\bigr].
 \label{eq:implemented-product-output-derivative}
\end{equation}
At $z_1=0$, this derivative equals $z_2$. Distinct lag coordinates can make
it arbitrarily large. Repeated occurrences of the same coordinate require
a separate calculation: $\tanh(z^2)$ has a bounded derivative.
Consequently, the conservative AST derivative test records what it can
verify; an infinite computed upper bound can represent overestimation.

The standard product-output mode checks a \emph{uniform second-moment
bound} by a separate project argument. Its growth envelope is
\begin{equation}
 \|F_{\mathcal T}(\mathbf{s})\|_2
 \leq B_{\mathrm{amp}}+
      \sum_{d=0}^{L-1} b_d^{\mathrm{amp}}\|\mathbf{s}_d\|_2,
 \qquad q_{\mathrm{amp}}:=\sum_d b_d^{\mathrm{amp}}<1.
 \label{eq:implemented-amplitude-envelope}
\end{equation}
Here $B_{\mathrm{amp}}$ covers constants and bounded nonlinear terms.
The nonnegative $b_d^{\mathrm{amp}}$ cover the magnitude of the remaining
unsaturated historical terms. They are \emph{growth coefficients}, whereas
$a_d^{\mathrm{GMC}}$ bound differences between histories.

To compute the envelope, a parameter contributes its absolute value to
the bounded part; a lag contributes one to its source/lag growth
coefficient; $\tanh(e)$ contributes a bounded part of one. Addition sums
these bounds. Multiplication by a bounded factor scales them. A product of
two unbounded envelopes is flagged for another construction. The mode
stacks the nonnegative growth coefficients into one target-by-source matrix
per lag and uses its largest singular value for $b_d^{\mathrm{amp}}$;
$B_{\mathrm{amp}}$ is the Euclidean norm of the componentwise bounded parts.
The mode
scales the complete drift using Eq.~\eqref{eq:implemented-gmc-scaling}
with the unscaled growth-coefficient sum in place of $\widetilde c$.
It stores the correspondingly scaled envelope for the executed equation.
Its root multiplier is stored as \texttt{AMPLITUDE\_SCALE}.

Let $M_0$ bound the $L^2$ norm of every initial lag vector. Minkowski's
inequality, the triangle inequality for root mean squared magnitudes, gives
\begin{equation}
 \|\mathbf x_{t+1}\|_{L^2}
 \leq B_{\mathrm{amp}}+
      \sum_d b_d^{\mathrm{amp}}\|\mathbf x_{t-d}\|_{L^2}
      +\sigma_\varepsilon\sqrt N.
 \label{eq:implemented-amplitude-recursion}
\end{equation}
Define
\begin{equation}
 M_{\mathrm{amp}}=
 \max\left\{M_0,\,
       \frac{B_{\mathrm{amp}}+\sigma_\varepsilon\sqrt N}
            {1-q_{\mathrm{amp}}}\right\}.
 \label{eq:implemented-amplitude-bound}
\end{equation}
If the historical vectors have $L^2$ norms at most $M_{\mathrm{amp}}$,
the next vector has norm at most
$B_{\mathrm{amp}}+q_{\mathrm{amp}}M_{\mathrm{amp}}
+\sigma_\varepsilon\sqrt N\leq M_{\mathrm{amp}}$.
Induction therefore yields
$\sup_{v\geq1}\mathbb E\|\mathbf x_{\tau+v}\|_2^2
\leq M_{\mathrm{amp}}^2$.
The uniform initial box gives the safe choice $M_0=r_0\sqrt N$.
The implementation stores these quantities in
\texttt{AmplitudeMoment\-Certificate}, with guarantee
\path{uniform_second_moment_bound}. This is our derived amplitude
criterion. Exponential forgetting and a stationary initialization require
their own additional conditions.

The pure saturated-product example
$x_{t+1}=\tanh(s_0s_1)+\varepsilon_{t+1}$ is particularly simple:
independent zero-mean Gaussian innovations give
$\mathbb E x_{t+1}^2\leq1+\sigma_\varepsilon^2$ at every generated step.
If two inputs are instead supplied by an already jointly stationary
external process, applying a fixed pointwise transformation preserves
their time-shift-invariant joint laws. In a feedback generator, the
history belongs to the newly defined recurrence, whose own growth or
contraction conditions must be checked as above.

For a scalar example
$F(\mathbf{s})=0.2s_0+0.4\tanh(s_0s_1)$, the growth coefficient is $0.2$
and the bounded part is $0.4$. The default multiplier is $0.95$, giving
$q_{\mathrm{amp}}=0.19$ and $B_{\mathrm{amp}}=0.38$.
With $r_0=1$ and $\sigma_\varepsilon=0.05$,
Eq.~\eqref{eq:implemented-amplitude-bound} gives $M_{\mathrm{amp}}=1$.
This example controls expected squared magnitude while allowing stochastic
fluctuations and sensitivity to the history.

\subsubsection{Initialization and Independent Innovations}
Both standard modes draw their initial $N\times L$ history directly from
$[-r_0,r_0]$, then generate the requested observations with independent
Gaussian innovations. The initial values are retained as a separate history
in the returned array; the subsequent entries are recursively generated.
A dataset window or diagnostic must select those generated entries explicitly.

The two standard transformations admit the configured sums/products through
their respective checks; a numerical verification failure is reported.
Legacy all-lag saturation and unsaturated sampling remain available for
reproducing earlier experiments. The unsaturated sampler rejects the entire
candidate mechanism when its conservative global derivative bound is
infinite, then samples a new skeleton, coefficients, and lags with a
deterministically derived attempt seed. This retry cost belongs to mechanism
sampling. A recorded zero-rejection rate for a standard mode describes the
chosen prior and execution.

Gaussian innovations have unbounded support. A finite amplitude threshold
can therefore be crossed even by a process with a uniform second-moment
bound or a stationary law. Tail flatness, low amplitude, and their observed
changes remain empirical diagnostics, as defined in
Section~\ref{sec:generator-diagnostics}.

\paragraph{Independent innovations and the interpretation of convergence.}
\label{sec:implemented-shao-wu-independent-noise}
The GMC comparison shares the actual numerical innovation at each future
step. Independent dataset trajectories receive distinct initialization and
innovation streams, derived from the root seed and trajectory index.

For the factor mode, the contraction and moment assumptions imply
forgetting of initialization and convergence in distribution to a common
stationary law. \citet[Theorem~2, p.~427; discussion pp.~426--428]{wu2004iterated}
give the underlying random-iteration result. Applying it separately to two
independent runs gives
\begin{equation}
 \operatorname{Law}(\mathbf x_{\tau+v})\Rightarrow\pi_{\mathrm{GMC}},
 \qquad
 \operatorname{Law}(\mathbf x'_{\tau+v})\Rightarrow\pi_{\mathrm{GMC}}.
 \label{eq:implemented-gmc-independent-laws}
\end{equation}
Independent trajectories can have different values while their laws
approach the same regime. For independent stationary copies, direct
expansion of the square yields
\begin{equation}
 \mathbb E\|\mathbf x_t-\mathbf x'_t\|_2^2
 =2\sum_{k=1}^{N}\operatorname{Var}_{\pi_{\mathrm{GMC}}}(x_{k,t}).
 \label{eq:implemented-gmc-independent-distance}
\end{equation}
The product-output amplitude record certifies the moment bound derived
above; distributional convergence is a separate question for that mode.
